\section{Selective positivity and calibrated sparse carry blocks}

The quadratic parent still has $31$ positive variables and $22$ equations.
A useful simplification is to shift only
\[
 F,\alpha,z,h,i,j,s,w,\eta,\zeta,\gamma,g
\]
by one, representing the other nineteen parent coordinates directly by
nonnegative integers. This reduces the symbolic coefficient records from
$136$ to $109$, and the distinct active monomials from $71$ to $62$
($39$ quadratic, $22$ linear, and the constant).

Here is the entire positivity argument. Write $q=1+\beta J$,
$C=q-F-\alpha$, $k=\eta+\zeta$, $Y=a-E$, $Q=c-\eta$, and $T=P-k$.
Transport gives $(K_0+w)C=F+z\beta J>0$, so $C>0$, $q\ge3$, and $J>0$.
The equations $L=q^2-1$, $N=q(L+1)$, and $X=wq$ make $L,N,X$ positive.
Next $Y=sN>0$, $E=XY>0$, and $a=E+Y>0$.
Then $Q=kY>0$, $c=Q+\eta>0$, $T=EQ>0$, $P=T+k>0$, and
$\tau^2=PT+1$ makes $\tau>0$.
The defining equations make $D,\Delta,C_2$ positive.
Since $I=iC_2>0$ and $B_{\rm aux}=I^2>0$,
$B_{\rm aux}=\Delta(F_2-1)$ and $F_2=f^2$ imply $f\ge2$.
Finally $y^2=B_{\rm aux}g+1>1$ and
$V^2=(B_{\rm aux}-1)g+1\ge1$ make $y,V$ positive;
$of=V+c>0$ makes $o>0$. Thus every unshifted coordinate is positive.
The inverse subtracts one from exactly the twelve guarded coordinates,
so this is a bijection of zero sets, not merely an existence reduction.
At the natural origin, residual~2 equals $-1$.

Let $p$ be the padding position, $w_i$ the distinct variable positions,
$W=\max(p,w_i)$, and $H_j$ the targets. For the independently certified
weighted source polynomial $D(t)$, put $d=\deg D$ and $K=d+2W$.
The positional isolation identity is
\[
 [t^{H_j}]\left(t^p+\sum_i x_it^{w_i}\right)^2D(t)=2r_j(x).
\]
Importantly, $d$, $\max H_j$, and $p$ are kept distinct.
For the coherent native tuple every weighted coefficient has magnitude
strictly below $M=2^{864013}$. Choose
$\kappa=864026$ and $a=2^\kappa=8192M$.
Set $B=ab^2$ and
\begin{align*}
 A&=\sum_jB^{H_j},&\ell&=\sum_iB^{w_i},&N_0&=B^{W+1},\\
 J_c&=2(B^p+g)^2D(B)+A,&
 S&=g+N_0J_c,&
 T&=N_0-1-(b-1)\ell+N_0(B-2)A.
\end{align*}
The sparse $+1$ target background permits signed lower coefficients:
incoming borrows are only $0$ or $-1$. At a target the normalized digit
avoids $B-2$ exactly when the selected coefficient is zero.
Consequently binary disjointness $S\mathbin{\&}T=0$ tests all residuals.
For $0\le g<4bB^W$, one has $0\le S,T<B^E$, where $E=K+W+2$.
These estimates cover interior padding because $B^p+g<5bB^W$.

For the nine-variable ratio bridge take
\[
 L=2bB^E,\qquad R=L(L+S+T+1)-(S+1),\qquad G_{\rm cap}=32bB^W.
\]
When $b=2^k$ and $B=2^v$, the mask has
$\operatorname{pop}(T)=(W+23)v-31k-22$ set bits.
The two-block complement identity therefore gives the exact binomial
threshold exponent $(K+2W+25)v-30k-20$.
For even $K$ this is the square of the polynomial
\[
 N=2^{15\kappa/2-10}B^{(K+2W+10)/2}.
\]
For odd $K$, replace this by
$N=2^{8\kappa-10}bB^{(K+2W+9)/2}$.
The endpoint $b=1$ introduces no zero, since its mask would force every
source digit to vanish. Thus positive $b$ can be an input directly.
For the $W=174$ packed certificate,
$(\deg R,\deg N,\deg G_{\rm cap})=(18386,4778,349)$.

The integer bridge permits an arbitrary polynomial threshold, so it is
better to retain $L=B^E$ and
\[
 Y=2^{\kappa(K+2W+25)-21}b^{,2(K+2W+25)-31}.
\]
This is not the nine-variable square $N^2$.
With $b=1+3f$, even $\kappa$ gives $b,B,L\equiv1\pmod3$,
$T\equiv-22\equiv2\pmod3$, and hence $X\equiv1+T\equiv0\pmod3$.
For completeness, both source coordinates $D$ and $\Delta$ are positive;
since their positions are distinct, at least one has positive weight.
Thus an encoded source zero satisfies $g\ge B\ge b$, even though the
natural source coordinates were not all assumed positive.

The source identities, coefficient bounds, signed-borrow cases and threshold
arithmetic are independently checked. The unbounded compiler/calibration and
Jones/Sun arithmetic bridges remain explicitly inherited mathematical proofs.
